% page 527 du fac-simile (image p-526 du scan)
% bloc 170.2 x 250.6 mm ; marge gauche 31.8 mm
\pgc{10.160mm}{0.000mm}{
\l{\cel{53}519.}
\l{}
\l{\sou{sep}. Sis plus un (6 +1, o 4 +3). - DEFIS.}
\l{}
\l{\sou{sepalo}. (\sou{bot.)}\cel{2}Folieto di la kalico di floro. - EFIS.}
\l{}
\l{\sou{separar.}\cel{1}(\sou{trans., de)}. I. Igar aparta, ici de iti, kozi, personi,}
\l{\cel{3}qui inter-kunesis (esis asemblita, unionita). - II. Esar situi-}
\l{\cel{3}ta inter kozi, personi, tale ke ici impedesas inter-kunesar. -}
\l{\cel{3}DEFIRS.}
\l{}
\l{\sou{sepio}. (\sou{zool.)}\cel{2}Molusko di qua la korpo esas karnoza e depresita}
\l{\cel{3}- qua havas sur sua dorso lamelo kalkoza, friabla (sepiosto)}
\l{\cel{3}e difuzas cirkum su, kande lu timas atakeso, liquido nigratra}
\l{\cel{3}quan onu uzas por facar farbi (sepio-koloro).- L.\cel{1}\sou{sepia}. DEFIRS}
\l{}
\l{\sou{septa}. (\sou{medic.}) Qua putrigas. - DEFIRS.}
\l{}
\l{\sou{septembro}. La monato nonesma di la yaro. - DEFIRS.}
\l{}
\l{\sou{septeto}. (\sou{muziko)}. Muziko koncertala por sep instrumenti o sep}
\l{\cel{3}voci (kun o sen akompano). - DEFRS.}
\l{}
\l{\sou{septiliono}. (\sou{matem.}) 1042}
\l{}
\l{\sou{septimo}. I. (skermo).Un ek la fazi di atako ; pareo di ta atako.}
\l{\cel{3}- II. (\sou{muzi}ko).Intervalo ek sep noti, de \textquotedbl{}do\textquotedbl{} til\textquotedbl{}si\textquotedbl{}, e qua}
\l{\cel{3}preiras la oktavo. -}
\l{}
\l{\sou{septua-}\cel{1}.Prefixo ciencala : \textquotedbl{}qua havas sep...\textquotedbl{}.}
\l{}
\l{\sou{septuagesimo}. (\sou{liturgio katolika}).La dio sepadekesma ante la}
\l{\cel{3}oktavo di pasko.}
\l{}
\l{\sou{sepultar}. (\sou{trans., en)}. Envelopar (la kadavro di persono) per}
\l{\cel{3}tuko, ante enterigar lu. - EFIS.}
\l{}
\l{\sou{sequar}. (\sou{trans.}) I. Irar dop (ulu, ulo) - (a) Irar dop (ulu),}
\l{\cel{3}akompanante ; (b) Esar dop o pos (ulo). - II. Direktar su}
\l{\cel{3}konforme ad (ulo, ulu) - (a) Durar, en spaco, segun la di-}
\l{\cel{3}reciono quan indikas ulu ; (\sou{metaf}.)egardar, pri sua konduto,}
\l{\cel{3}la precepto, la modereso-grado quan indikis ulu, ulo ; (b)}
\l{\cel{3}durar irar segun la direciono,kondutar segun sua selekto.}
\l{\cel{3}- EFIS.}
\l{}
\l{sequencio. (liturgio katol.) . Versajo ek versi segun-mezuro e}
\l{\cel{3}ritmoza qua, en la mesi solena, sequas la graduelo (fragmento}
\l{\cel{3}di la oficio, inter la epistolo e la prozo, ante Evangelio,}
\l{\cel{3}quan onu pronuncis olim sur la gradi di la jubeo o di la ambono)}
\l{\cel{4}e la \textquotedbl{}aleluya\textquotedbl{}. -}
\l{}
\l{sequestrar. (trans.) (yurocienco).Depozar (la kozo di la litijo-}
\l{\cel{3}punto) en la manui di altra persono qua gardos lu til kande}
\l{\cel{3}decidesabos pri la persono a qua lu devas apartenar. DEFIRS.}
\l{}
\l{\cel{1}\sou{sequoyo}. (\sou{bot.)}\cel{2}Genero de koniferi taxodinea, qua kreskas en Ka-}
\l{\cel{3}lifornia (Usa) - arboro imensa di qua ula sorto esas alta ye}
\l{\cel{4}140 metri. - L.\cel{1}\sou{sequoja}. - EF.}
}
